% -----------------------------
% Interlineado
% -----------------------------
\usepackage{setspace}
\onehalfspacing

% -----------------------------
% Idioma y codificación
% -----------------------------
\usepackage[spanish]{babel}
\usepackage{csquotes} % Recomendado para evitar advertencias con babel y biblatex
\usepackage{caption} %esta cosa es importante porque resuleve el problema del idioma español y babel
% -----------------------------
% Tipografía (Calibri)
% -----------------------------
\usepackage{fontspec}
\setmainfont{Calibri}

% -----------------------------
% Márgenes (Estilo Anteproyecto)
% -----------------------------
\usepackage{geometry}
\geometry{
	letterpaper,
	left=2.5cm,
	right=2.5cm,
	top=2.5cm,
	bottom=2.5cm
}

% -----------------------------
% Imágenes
% -----------------------------
\usepackage{graphicx} % Permite insertar imágenes	

% -----------------------------
% Párrafos (Sin sangría, espacio entre párrafos)
% -----------------------------
\setlength{\parindent}{0pt}
\setlength{\parskip}{0.8em}

% -----------------------------
% Formato de Títulos
% -----------------------------
\usepackage{titlesec}

% --- Numeración de secciones SIN capítulo ---
\renewcommand{\thesection}{\arabic{section}}

% --- Formato del título ---
\titleformat{\section}
{\normalfont\bfseries}
{\thesection.}
{1em}
{}

% --- Sangría del título ---
\titlespacing*{\section}
{1.5em}  % sangría a la derecha
{2.5ex}  % espacio antes
{1.5ex}  % espacio después

% -----------------------------
% Bibliografía (Formato APA)
% -----------------------------
\usepackage[
backend=biber, 
	style=apa % Cambiado a APA nativo para biblatex
	]{biblatex}
\addbibresource{bibliografia/referencias.bib}

% -----------------------------
% Hipervínculos
% -----------------------------
\usepackage{hyperref}
\hypersetup{
	colorlinks=true,
	linkcolor=black,
	urlcolor=black,
	citecolor=black
}

% -----------------------------
% Profundidad del Índice y Numeración
% -----------------------------
% El número 3 le dice a LaTeX que llegue hasta \subsubsection (nivel 3)
\setcounter{secnumdepth}{3} % Numera hasta sub-subsección (ej. 1.1.1)
\setcounter{tocdepth}{3}    % Muestra hasta sub-subsección en el índice

% -----------------------------
% Formato y Tamaños de Títulos
% -----------------------------
\usepackage{titlesec}
\renewcommand{\thesection}{\arabic{section}} % Hace que empiece en 1 y no en 0.1

% 1. Sección: Letra grande (\Large)
\titleformat{\section}{\normalfont\Large\bfseries}{\thesection.}{1em}{}

% 2. Subsección: Letra mediana (\large)
\titleformat{\subsection}{\normalfont\large\bfseries}{\thesubsection.}{1em}{}

% 3. Sub-subsección: Letra normal (\normalsize)
\titleformat{\subsubsection}{\normalfont\normalsize\bfseries}{\thesubsubsection.}{1em}{}